% Gerado por scripts/migrate_markdown_to_latex.py.
% Fonte migrada: docs/apresentacoes/simposio_2026/guias/guia_estudo.md


\chapter{Guia de Estudo do Tema}



\section{1. O que é um Transformer}





Transformer é uma arquitetura de rede neural introduzida por Vaswani et al. em 2017. Ela se tornou central em IA porque substituiu estruturas recorrentes por operações matriciais paralelizáveis. Isso permitiu treinar e executar modelos muito maiores com melhor uso de hardware moderno.




Em termos simples, o Transformer recebe uma sequência de vetores e aprende quais partes dessa sequência devem influenciar outras partes. Ele faz isso principalmente por meio do mecanismo de atenção.



\section{2. Por que a atenção custa caro}





Na autoatenção, cada token pode ser comparado com muitos outros tokens da mesma sequência. Isso significa que, quando a sequência cresce, o número de comparações cresce bastante. Além disso, antes da atenção, o modelo ainda faz projeções lineares densas para gerar Q, K e V.



Então, o custo da inferência não vem de um único ponto, mas da combinação de:

\begin{itemize}
\item multiplicações matriciais densas
\item comparação entre tokens
\item armazenamento e leitura de memória intermediária
\end{itemize}


\section{3. Diferença entre treino e inferência}




Treino é a fase em que o modelo aprende ajustando seus pesos. Inferência é a fase em que o modelo já treinado é usado para gerar respostas.




No seu trabalho, o foco está na inferência. Isso é importante porque muitas otimizações relevantes para uso prático buscam tornar a inferência mais barata, rápida e escalável.



\section{4. O que é quantização}




Quantização é a redução da precisão numérica usada para representar dados. Por exemplo, em vez de armazenar valores em formatos de maior precisão, é possível usar 8 bits ou menos.



Benefícios esperados:

\begin{itemize}
\item menos memória
\item menos tráfego de dados
\item potencial aumento de velocidade
\end{itemize}


Risco:

\begin{itemize}
\item introdução de erro numérico
\end{itemize}


Por isso, quantização sempre envolve um trade-off entre eficiência e fidelidade.



\section{5. O que é pruning}




Pruning é a remoção ou desativação de pesos considerados menos relevantes. Em tese, isso reduz o tamanho efetivo do modelo e o número de operações.



Mas existe uma diferença importante em relação à quantização:

\begin{itemize}
\item quantização muda a representação dos valores
\item pruning muda a estrutura efetiva do cálculo
\end{itemize}


Na prática, pruning pode exigir mais cuidado para realmente virar ganho de execução, porque hardware e bibliotecas nem sempre aproveitam bem esparsidade.



\section{6. O que é KV cache}




KV cache significa key-value cache. Durante a geração de tokens em modelos autoregressivos, o modelo guarda chaves e valores já calculados para não recomputar tudo a cada passo.



Isso acelera a inferência, mas tem um custo:

\begin{itemize}
\item conforme o contexto cresce, o KV cache ocupa cada vez mais memória
\end{itemize}


Por isso, técnicas que comprimem KV cache são muito relevantes para inferência eficiente.



\section{7. O que o TurboQuant propõe}




TurboQuant é um método recente de quantização online divulgado pelo Google Research em 24 de março de 2026. O paper correspondente está no arXiv desde 28 de abril de 2025.




Em alto nível, ele busca quantizar vetores com baixa distorção, incluindo aplicações em KV cache e busca vetorial. Na narrativa da sua apresentação, o ponto principal não é reproduzir toda a matemática do método, mas usar o trabalho como evidência de que quantização está no centro das soluções mais fortes para eficiência em inferência.



\section{8. Por que quantização parece mais promissora aqui}



No contexto do seu trabalho, quantização parece uma primeira hipótese mais forte por três motivos:


\begin{enumerate}
\item Ataca diretamente memória e largura de banda.
\item Tem apoio forte em literatura recente.
\item Pode ser comparada de forma relativamente controlada com baseline e pruning.
\end{enumerate}



Isso não significa que pruning é irrelevante. Significa apenas que, para um primeiro ciclo experimental, quantização parece oferecer melhor custo-benefício científico.



\section{9. O que significam as métricas}



\subsection{Latência}


Tempo gasto para executar uma operação ou cenário.



\subsection{Memória}


Quanto espaço é necessário durante a execução.



\subsection{FLOPs}


Estimativa do número de operações aritméticas de ponto flutuante.



\subsection{MSE}


Erro quadrático médio. Penaliza mais fortemente erros grandes.



\subsection{MAE}


Erro absoluto médio. Mede, em média, o quanto a saída se distancia da referência.



\subsection{R²}


Coeficiente de determinação. Indica o quanto a saída candidata acompanha a referência.



\subsection{Similaridade de cosseno}



Mede o alinhamento entre dois vetores. É útil quando interessa preservar direção e estrutura relativa da saída.



\section{10. Como explicar seu experimento de forma clara}



Uma forma simples de explicar é:



'Eu fixo os mesmos dados e os mesmos pesos-base, executo operações centrais do Transformer em três cenários diferentes e comparo custo computacional e qualidade de saída.'



Essa frase já resume:

\begin{itemize}
\item controle experimental
\item comparação entre cenários
\item foco em eficiência e fidelidade
\end{itemize}


\section{11. Perguntas prováveis da banca}



\subsection{'Por que você escolheu Transformer?'}


Porque é a arquitetura dominante em IA atual e concentra desafios reais de inferência em tempo e memória.



\subsection{'Por que comparar quantização com pruning?'}


Porque são duas estratégias clássicas e relevantes de otimização, mas atuam de formas diferentes: uma na representação numérica e outra na estrutura efetiva do modelo.



\subsection{'Por que dizer que quantização é melhor se você ainda não rodou o experimento?'}


Eu não afirmo como conclusão do meu experimento. Eu apresento como hipótese inicial motivada por literatura recente e por características do problema.



\subsection{'Por que usar cenário controlado e modelo simplificado?'}


Porque isso ajuda a isolar o efeito da técnica de otimização antes de ir para modelos completos mais complexos.



\subsection{'O que você espera encontrar?'}


Espero observar redução de custo computacional nos cenários otimizados e depois comparar o quanto isso preserva ou degrada a qualidade das saídas em relação ao baseline.



\section{12. Frase de segurança para a apresentação}



Se em algum momento você travar, pode usar esta frase:



'Meu trabalho ainda está na fase inicial de execução, então hoje eu apresento principalmente a relevância do problema, a hipótese técnica e o desenho experimental que vai permitir validar essa hipótese com dados próprios.'
